\chapter{Implementation and Evaluation}

\section{Project Source-Code Structure}
    The project is organised into multiple separate files by module, corresponding to the group's Git branching model. Each class is split into a header file (\texttt{.h}) and an implementation file (\texttt{.cpp}), which reduces compilation dependencies and allows members to work in parallel. The directory structure and division of work are presented in Table~\ref{tab:phancong} at the end of this chapter.

\section{Implementation of the Core Entities Module}
    \subsection{The Abstract Account Class}
    Listing~\ref{lst:account} illustrates the declaration of the abstract class \texttt{Account}. The \texttt{protected} attributes, the virtual destructor, the two pure virtual functions \texttt{printDetails()} and \texttt{clone()}, together with the business methods, can be clearly seen.

\begin{lstlisting}[caption={Declaration of the abstract Account class}, label={lst:account}]
class Account {                       // <<abstract>>
protected:                            // protected for derived classes
    std::string   accountID;
    double        balance;
    std::string   ownerName;
    AccountStatus status;
public:
    Account(std::string id, std::string owner, double bal);
    virtual ~Account() = default;     // virtual destructor

    std::string   getAccountID() const;
    double        getBalance()   const;
    AccountStatus getStatus()    const;

    // realistic banking operations
    void freeze();
    void unfreeze();
    bool withdraw(double amt);

    virtual void     printDetails() const = 0;   // pure virtual
    virtual Account* clone()        const = 0;    // Virtual Constructor
};
\end{lstlisting}

    \subsection{Multiple Inheritance in InternationalTransfer}
    Listing~\ref{lst:intl} illustrates the multiple-inheritance syntax. The \texttt{InternationalTransfer} class inherits from both \texttt{Transaction} and \texttt{ForeignExchange}, and calls the constructors of both base classes in its initialiser list.

\begin{lstlisting}[caption={Multiple inheritance in InternationalTransfer}, label={lst:intl}]
// Multiple inheritance: Transaction + ForeignExchange
// The two bases share no ancestor -> NO virtual inheritance needed
class InternationalTransfer : public Transaction, public ForeignExchange {
    std::string swiftCode;
public:
    InternationalTransfer(std::string id, double amt, std::string from,
                          std::string to, Date d, Location loc,
                          std::string src, std::string tgt,
                          double rate, std::string swift)
        : Transaction(id, amt, from, to, d, loc),
          ForeignExchange(src, tgt, rate),
          swiftCode(swift) {}
    std::string  getDetails() const override;
    Transaction* clone()      const override;
};
\end{lstlisting}

\section{Implementation of the Audit Rules Module -- Polymorphism}
    \subsection{The AuditRule Interface}
    Listing~\ref{lst:auditrule} shows the abstract interface for all rules. Note that the \texttt{evaluate()} method receives a \texttt{const DataManagement\&} reference so that rules requiring historical context can access it.

\begin{lstlisting}[caption={The abstract AuditRule interface}, label={lst:auditrule}]
class AuditRule {                     // <<interface>>
public:
    virtual ~AuditRule() = default;
    // Receives the DB so pattern-based rules can access history
    virtual AuditInspectionResult evaluate(
        const Transaction* t,
        const DataManagement& db) const = 0;
    virtual std::string getRuleName() const = 0;
    virtual Severity    getSeverity() const = 0;
};
\end{lstlisting}

    \subsection{Implementation of the Smurfing Detection Rule}
    Listing~\ref{lst:smurfing} shows how \texttt{SmurfingRule} traverses the entire transaction history to count small transactions on the same day from the same sender -- which is only possible thanks to the \texttt{db} parameter.

\begin{lstlisting}[caption={Smurfing detection logic}, label={lst:smurfing}]
AuditInspectionResult SmurfingRule::evaluate(
        const Transaction* t, const DataManagement& db) const {
    AuditInspectionResult r; r.transactionID = t->getID();
    if (t->getAmount() >= smallTxThreshold) return r; // only small tx
    int count = 0;
    for (auto* o : db.getTransactions())
        if (o->getSenderID() == t->getSenderID()
            && o->getDate().sameDay(t->getDate())
            && o->getAmount() < smallTxThreshold) ++count;
    if (count >= maxFrequencyPerDay) {
        r.isPassed = false;
        r.ruleViolated = getRuleName();
        r.severity = getSeverity();
    }
    return r;
}
\end{lstlisting}

    \subsection{The Polymorphic Audit Loop}
    The heart of the system lies in the \texttt{runAudit()} method of \texttt{AuditEngine} (Listing~\ref{lst:runaudit}). The call \texttt{rule->evaluate(t, db)} is exactly where the dynamic dispatch mechanism takes effect: the same line of code executes different logic depending on the actual rule type.

\begin{lstlisting}[caption={The audit loop with dynamic dispatch}, label={lst:runaudit}]
std::vector<AuditInspectionResult>
AuditEngine::runAudit(const DataManagement& db) const {
    std::vector<AuditInspectionResult> out;
    for (auto* t : db.getTransactions())
        for (auto* rule : rules) {
            // DYNAMIC DISPATCH: choose the right rule via vtable
            AuditInspectionResult res = rule->evaluate(t, db);
            if (!res.isPassed) out.push_back(res);
        }
    return out;
}
\end{lstlisting}

\section{Memory Management -- The Rule of Five}
    Because \texttt{DataManagement} owns dynamically allocated pointers, it must define five special functions for safe memory management: the destructor, the copy constructor, the copy assignment operator, the move constructor, and the move assignment operator. The challenge is that when copying a \texttt{vector<Account*>}, copying only the pointers would make two objects point to the same memory region (a \textit{shallow copy}), leading to a double-free error.

    The solution is the \textbf{Virtual Constructor} idiom \cite{Gamma1994}: each derived class implements the virtual method \texttt{clone()} that returns a copy of the correct type of itself. Thanks to this, the copy constructor can perform a \textbf{polymorphic deep copy} as shown in Listing~\ref{lst:ruleof5}.

\begin{lstlisting}[caption={Polymorphic deep copy via clone()}, label={lst:ruleof5}]
void DataManagement::deepCopyFrom(const DataManagement& o) {
    for (auto* a : o.accounts)
        accounts.push_back(a->clone());      // copy as the derived type
    for (auto* t : o.transactions)
        transactions.push_back(t->clone());
}
DataManagement::~DataManagement() { freeAll(); }   // free everything
\end{lstlisting}

\section{Experimental Results}
    \subsection{Interactive Interface}
    The system provides an interactive menu that allows the user to enter data, view lists, select audit rules, and export reports. The menu also includes an integrated Help \& Guide option that explains the input format directly within the program. Figure~\ref{fig:menu} shows the main menu and the report-export operation.

    \begin{figure}[H]
        \centering
        \includegraphics[width=\textwidth]{menu.png}
        \caption{The interactive menu of the system and the report-export operation.}
        \label{fig:menu}
    \end{figure}

    \subsection{Running the Audit}
    After loading the sample data set of 4 accounts and 8 transactions, the user can select any combination of audit rules to run. Figure~\ref{fig:audit_output} illustrates the case where the \texttt{GeoAnomalyRule} is selected: the system correctly detects transaction \texttt{TX007} originating from Iran -- a high-risk country -- and flags it with a \texttt{MEDIUM} severity level. Each rule can be enabled independently, and selecting all four rules yields the full set of warnings.

    \begin{figure}[H]
        \centering
        \includegraphics[width=\textwidth]{audit_output.png}
        \caption{Selecting an audit rule and the corresponding detection result.}
        \label{fig:audit_output}
    \end{figure}

    \subsection{Exported Report File}
    The audit result is also exported to a plain-text file \texttt{audit\_report.txt}, fulfilling the requirement of storing data in files. Figure~\ref{fig:report_output} shows the content of the exported report, including the total number of warnings and a detailed description of each violation.

    \begin{figure}[H]
        \centering
        \includegraphics[width=\textwidth]{report_output.png}
        \caption{Content of the exported audit report file.}
        \label{fig:report_output}
    \end{figure}

    \subsection{Memory Management Verification}
    To prove that the program does not leak memory, the group recompiled it with the \texttt{-fsanitize=address,leak} flag and ran it. The AddressSanitizer tool confirmed that no memory leaks were detected, validating the correctness of the Rule of Five and the Virtual Constructor idiom.

\section{Division of Work}
    The group applies a feature-based Git branching model. The four members take charge of the four architectural modules; each member ``owns'' one OOP principle to present, ensuring no overlap and that every member participates in coding. The details are presented in Table~\ref{tab:phancong}.

    \begin{table}[H]
        \centering
        \caption{Division of work and responsibilities of members}
        \label{tab:phancong}
        \begin{tabular}{|l|p{4cm}|p{3.8cm}|}
            \hline
            \textbf{Member} & \textbf{Files in charge} & \textbf{Key principle} \\ \hline
            Mai Nam Khanh \textit{(Leader)} & \texttt{DataManagement}, \texttt{AuditEngine}, \texttt{FileManager}, \texttt{ConsoleUI} & Memory management, Aggregation, I/O \\ \hline
            Hoang Le Anh Duc & \texttt{Account.h/.cpp} & Inheritance, Encapsulation \\ \hline
            Vu Duy Tuan & \texttt{Transaction.h/.cpp}, \texttt{Common.h} & Multiple inheritance, Composition \\ \hline
            Pham Duc Tri & \texttt{AuditRule.h/.cpp} & Abstraction, Polymorphism \\ \hline
        \end{tabular}
    \end{table}
